\documentclass{article}

% packages
\usepackage{amsmath, amsthm, thmtools, amsfonts, amssymb, luacode, catchfile, hyperref, ifthen, graphicx, bbm}
\ifcsname c@kobocompile\endcsname
	\usepackage[a5paper, total={1072pt, 1448pt}, margin=10pt, includeheadfoot]{geometry} % set page margins
\else
	\usepackage[a4paper, margin=50pt, includeheadfoot]{geometry}
\fi
\usepackage[shortlabels]{enumitem}
\usepackage[skip=3pt, indent=0pt]{parskip}

% language
\usepackage[bidi=basic, layout=tabular, provide=*]{babel}
\ifcsname c@english\endcsname
	\babelprovide[main, import]{english}
\else
	\babelprovide[main, import, alph=letters.gershayim]{hebrew}
	\setlist[enumerate,2]{label={\alph*.}}
	\babelprovide{rl}
\fi
\babelfont{rm}[Renderer=Harfbuzz]{Libertinus Serif}
\babelfont{sf}{Libertinus Sans}
\babelfont{tt}{Libertinus Mono}

% style
\ifcsname c@stoppagebreaks\endcsname
\else
	\AddToHook{cmd/section/before}{\clearpage}	% Add line break before section
\fi
\linespread{1.3}
\setcounter{secnumdepth}{0}		% Remove default number tags from sections, this won't do well with theorems
\AtBeginDocument{\setlength{\belowdisplayskip}{3pt}}
\AtBeginDocument{\setlength{\abovedisplayskip}{3pt}}
\graphicspath{ {../images/} }

% operators
\DeclareMathOperator\cis{cis}
\DeclareMathOperator\Sp{Span}
\DeclareMathOperator\tr{tr}
\DeclareMathOperator\im{Im}
\DeclareMathOperator\re{Re}
\DeclareMathOperator\diag{diag}
\DeclareMathOperator*\lowlim{\underline{lim}}
\DeclareMathOperator*\uplim{\overline{lim}}
\DeclareMathOperator\rng{rng}
\DeclareMathOperator\Sym{Sym}
\DeclareMathOperator\Arg{Arg}
\DeclareMathOperator\Log{Log}
\DeclareMathOperator\dom{dom}
\DeclareMathOperator\supp{Supp}
\DeclareMathOperator\var{Var}
\DeclareMathOperator\cov{Cov}
\DeclareMathOperator\End{End}
\DeclareMathOperator\id{id}
\DeclareMathOperator\cl{cl}
\DeclareMathOperator\gal{Gal}
\DeclareMathOperator\dist{dist}
\DeclareMathOperator\form{form}
\DeclareMathOperator\term{term}
\DeclareMathOperator\sent{sent}
\DeclareMathOperator\Th{Th}
\DeclareMathOperator\Hom{Hom}
\DeclareMathOperator\Aut{Aut}
\DeclareMathOperator\Char{char}
\DeclareMathOperator\vol{vol}
\DeclareMathOperator\GL{GL}
\DeclareMathOperator\SL{SL}
\DeclareMathOperator\lcm{lcm}

% commands
\newcommand{\cupdot}{\mathbin{\mathaccent\cdot\cup}}
\newcommand{\indicator}{\mathbbm{1}}
\newcommand{\AAA}[0]{\mathbb{A}}
\newcommand{\Aa}[0]{\mathcal{A}}
\newcommand{\BB}[0]{\mathbb{B}}
\newcommand{\Bb}[0]{\mathcal{B}}
\newcommand{\CC}[0]{\mathbb{C}}
\newcommand{\Cc}[0]{\mathcal{C}}
\newcommand{\DD}[0]{\mathbb{D}}
\newcommand{\Dd}[0]{\mathcal{D}}
\newcommand{\EE}[0]{\mathbb{E}}
\newcommand{\FF}[0]{\mathbb{F}}
\newcommand{\Ff}[0]{\mathcal{F}}
\newcommand{\Ii}[0]{\mathcal{I}}
\newcommand{\Gg}[0]{\mathcal{G}}
\newcommand{\HH}[0]{\mathbb{H}}
\newcommand{\Hh}[0]{\mathcal{H}}
\newcommand{\Kk}[0]{\mathcal{K}}
\newcommand{\Ll}[0]{\mathcal{L}}
\newcommand{\Mm}[0]{\mathcal{M}}
\newcommand{\NN}[0]{\mathbb{N}}
\newcommand{\Nn}[0]{\mathcal{N}}
\newcommand{\PP}[0]{\mathbb{P}}
\newcommand{\Pp}[0]{\mathcal{P}}
\newcommand{\QQ}[0]{\mathbb{Q}}
\newcommand{\RR}[0]{\mathbb{R}}
\newcommand{\Rr}[0]{\mathcal{R}}
\newcommand{\Ss}[0]{\mathcal{S}}
\newcommand{\TT}[0]{\mathbb{T}}
\newcommand{\Tt}[0]{\mathcal{T}}
\newcommand{\UU}[0]{\mathbb{U}}
\newcommand{\Uu}[0]{\mathcal{U}}
\newcommand{\Vv}[0]{\mathcal{V}}
\newcommand{\Ww}[0]{\mathcal{W}}
\newcommand{\Xx}[0]{\mathcal{X}}
\newcommand{\Yy}[0]{\mathcal{Y}}
\newcommand{\ZZ}[0]{\mathbb{Z}}
\newcommand{\Zz}[0]{\mathcal{Z}}
\newcommand{\acts}[0]{\circlearrowright}
\newcommand{\explain}[2] {
	\begin{flalign*}
		 && \text{#2} && \text{#1}
	\end{flalign*}
}
\newcommand{\maketitleprint}[1][green]{
	\begin{center}
		\includegraphics[width=6cm]{dragon_#1}
	\end{center}
}

% theorem commands
\newtheoremstyle{c_remark}
	{}	% Space above
	{}	% Space below
	{}% Body font
	{}	% Indent amount
	{\bfseries}	% Theorem head font
	{}	% Punctuation after theorem head
	{.5em}	% Space after theorem head
	{\thmname{#1}\thmnumber{ #2}\thmnote{ \normalfont{\text{(#3)}}}}	% head content
\newtheoremstyle{c_definition}
	{3pt}	% Space above
	{3pt}	% Space below
	{}% Body font
	{}	% Indent amount
	{\bfseries}	% Theorem head font
	{}	% Punctuation after theorem head
	{.5em}	% Space after theorem head
	{\thmname{#1}\thmnumber{ #2}\thmnote{ \normalfont{\text{(#3)}}}}	% head content
\newtheoremstyle{c_plain}
	{3pt}	% Space above
	{3pt}	% Space below
	{\itshape}% Body font
	{}	% Indent amount
	{\bfseries}	% Theorem head font
	{}	% Punctuation after theorem head
	{.5em}	% Space after theorem head
	{\thmname{#1}\thmnumber{ #2}\thmnote{ \text{(#3)}}}	% head content

\ifcsname c@english\endcsname
	\theoremstyle{plain}
	\newtheorem{theorem}{Theorem}[section]
	\newtheorem*{theorem*}{Theorem}
	\newtheorem{lemma}[theorem]{Lemma}
	\newtheorem{proposition}[theorem]{Proposition}
	\newtheorem*{proposition*}{Proposition}

	\theoremstyle{definition}
	\newtheorem{definition}[theorem]{Definition}
	\newtheorem*{definition*}{Definition}
	\newtheorem{example}{Example}[section]
	\newtheorem{exercise}{Exercise}[section]

	\theoremstyle{remark}
	\newtheorem*{remark}{Remark}
	\newtheorem*{solution}{Solution}
	\newtheorem{corollary}[theorem]{Corollary}
	\newtheorem{notation}[theorem]{Notation}
\else
	\theoremstyle{c_plain}
	\newtheorem{theorem}{משפט}[section]
	\newtheorem*{theorem*}{משפט}
	\newtheorem{lemma}[theorem]{למה}
	\newtheorem{proposition}[theorem]{טענה}
	\newtheorem*{proposition*}{טענה}

	\theoremstyle{c_definition}
	\newtheorem{definition}[theorem]{הגדרה}
	\newtheorem*{definition*}{הגדרה}
	\newtheorem{example}{דוגמה}[section]
	\newtheorem{exercise}{תרגיל}[section]

	\theoremstyle{c_remark}
	\newtheorem*{remark}{הערה}
	\newtheorem*{solution}{פתרון}
	\newtheorem{corollary}[theorem]{מסקנה}
	\newtheorem{notation}[theorem]{סימון}
\fi

% Questions related commands
\newcounter{question}
\setcounter{question}{1}
\newcounter{sub_question}
\setcounter{sub_question}{1}

\ifcsname c@english\endcsname
	\newcommand{\question}[1][0]{
		\ifthenelse{#1 = 0}{}{\setcounter{question}{#1}}
		\section{Question \arabic{question}}
		\addtocounter{question}{1}
		\setcounter{sub_question}{1}
	}

	\newcommand{\subquestion}[1][0]{
		\ifthenelse{#1 = 0}{}{\setcounter{sub_question}{#1}}
		\subsection{Part \alph{sub_question}}
		\addtocounter{sub_question}{1}
	}
\else
	\newcommand{\question}[1][0]{
		\ifthenelse{#1 = 0}{}{\setcounter{question}{#1}}
		\section{שאלה \arabic{question}}
		\addtocounter{question}{1}
		\setcounter{sub_question}{1}
	}

	\newcommand{\subquestion}[1][0]{
		\ifthenelse{#1 = 0}{}{\setcounter{sub_question}{#1}}
		\subsection{סעיף \localecounter{letters.gershayim}{sub_question}}
		\addtocounter{sub_question}{1}
	}
\fi

% import lua and start of document
\directlua{common = require ('../common')}

\GetEnv{AUTHOR}

% headers
\author{\AUTHOR}
\date\today

\title{פתרון מטלה 02 --- יסודות באלגברה קומוטטיבית}

\DeclareMathOperator{\Nil}{Nil}

\begin{document}
\maketitle
\maketitleprint[red]

\question{}
יהי $R$ חוג, נראה כי אין איברים נילפוטנטיים ב־$R / \Nil(R)$ פרט ל־$0$.
\begin{proof}
	יהי $x \in R$.
	אם $x \in \Nil(R)$ אז קיים $n \in \NN$ כך ש־$x^n = 0$ ולכן גם ${(x + \Nil(R))}^n = x^n + \Nil(R) = \bar{0}$, ונסיק ש־$\bar{0}$ הוא נילפוטנטי בחבורת המנה.
	נניח אם כך ש־$x \notin \Nil(R)$, ולכן לכל $n \in \NN$ מתקיים $x^n \ne 0$, ולכן גם $x^n + \Nil(R) = \Nil(R)$, כלומר $x + \Nil(R)$ לא נילפוטנטי בחוג המנה.
\end{proof}

\question{}
יהי $x \in R$ נילפוטנטי ונראה ש־$1 + x$ הפיך.
נסיק שסכום של הפיך ונילפוטנטי הוא הפיך.
\begin{proof}
	נסמן $n \in \NN$ המקיים $x^n = 0$ וניזכר בזהות,
	\[
		k^n - 1
		= (1 + \cdots + k^{n - 1})(k - 1)
	\]
	הנכונה לכל $k$, בפרט אם נציב $k = -x$ נקבל,
	\[
		{(-x)}^n - 1
		= \left(\sum_{i = 0}^{n - 1} {(-1)}^i x^i\right) \cdot ((-x) - 1)
	\]
	אבל $x^n = 0$ ולכן גם $-x^n = 0$ ובכל מקרה,
	\[
		0 - 1
		= -\left(\sum_{i = 0}^{n - 1} {(-1)}^i x^i\right) \cdot (x + 1)
	\]
	ולכן מצאנו ש־$x + 1$ הפיך וההופכי שלו מתקבל על־ידי $\sum_{i = 0}^{n - 1} {(-1)}^i x^i$.

	יהי $y \in R$ הפיך כלשהו, אז $\frac{x}{y} \in R$ איבר נילפוטנטי שכן ${(\frac{x}{y})}^n = \frac{x^n}{y^n} = \frac{0}{y^n} = 0$.
	אם כך אז $y + x = y (1 + \frac{x}{y})$, אבל האחרון הפיך כפי שראינו, ומכפלת הפיכים היא הפיכה.
\end{proof}

\question{}
יהי $R$ חוג, נראה כי $R \setminus R^\times$ אידאל אם ורק אם יש ל־$R$ אידאל מקסימלי יחיד, כלומר ש־$R$ הוא חוג מקומי.
\begin{proof}
	נניח ש־$R \setminus R^\times$ הוא אידאל, נראה שהוא מקסימלי ויחיד.
	לכל אידאל $R \setminus R^\times \subsetneq I \subset R$ בהכרח יש איבר $x \in R^\times$ כך ש־$x \in I$, לכן גם $1 \in I$, נסיק ש־$R \setminus R^\times$ אכן אידאל מקסימלי.
	ברור כי הוא יחיד מאותו שיקול של הכלת איברים הפיכים וסגירות לכפל.

	נניח מהצד השני ש־$I \subset R$ הוא אידאל מקסימלי יחיד ונראה ש־$I = R \setminus R^\times$, בפרט שזהו אידאל.
	משיקולי הפיכות ברור ש־$I \subset R \setminus R^\times$, לכן נותר להראות שאם $x \in R \setminus R^\times$ אז $x \in I$.
	על־ידי שימוש בעובדה ש־$(x) \subsetneq R$ ושימוש בלמה של צורן נוכל להסיק שקיים אידאל מקסימלי שכולל את $x$, לכן מיחידות בהכרח $x \in I$ כפי שרצינו.
\end{proof}

\question{}
יהי $R$ חוג ונראה שהתנאים הבאים שקולים:
\begin{enumerate}
	\item ל־$R$ יש אידאל ראשוני יחיד
	\item כל איבר ב־$R$ הוא הפיך או נילפוטנטי
	\item $R / \Nil(R)$ שדה
\end{enumerate}
\begin{proof}
	$1 \implies 2$:
	יהי $I$ אידאל מקסימלי כלשהו של $R$, לכן הוא אידאל ראשוני ולכן יחיד, כלומר ישנו אידאל מקסימלי יחיד לחוג.
	משאלה 3 נובע ש־$I = R \setminus R^\times$, בפרט אם $x \notin I$ אז הוא הפיך.
	יהי $x \in I$, מיחידות האידאל המקסימלי נקבל ש־$(x) \subset I$, ולכן $x$ נמצא בחיתוך כל האידאלים הראשוניים ולכן נילפוטנטי.

	$2 \iff 3$:
	ברור כנתון שבחוג המנה כל האיברים הם או הפיכים או מזדהים עם $0$.

	$2 + 3 \implies 1$:
	נתון כי $\Nil(R)$ הוא מקסימלי ולכן ראשוני ונראה שכל אידאל ראשוני שווה בדיוק לו.
	יהי $I$ אידאל ראשוני כלשהו ולכן ניתן להרחבה לאידאל מקסימלי, אך יש רק אחד כזה ולכן $I \subset \Nil(R)$.
	נניח ש־$x \in \Nil(R)$, אז נבחר $y \in I$ נילפוטנטי כך שהחזקה המאפסת שלו שונה משל $x$.
	\[
		{(x + y)}^n
		= 0 + x^{n - 1} + \cdots + y^n
	\]
	עבור $n$ כך ש־$x^n = 0$, אבל מסגירות לכפל וחיבור נוכל להסיק ש־${(x + y)}^n \in I$ ולכן מראשוניות גם $x + y \in I$ וכן $x \in I$.
\end{proof}

\question{}
יהי $R$ חוג כך שלכל אידאל $I \triangleleft R$ כך ש־$I \not\subset \Nil(R)$ קיים $0 \ne x \in I$ כך ש־$x^2 = x$.
נראה ש־$\Nil(R) = \Jj(R)$.
\begin{proof}
	מספיק להראות ש־$\Jj(R) \subset \Nil(R)$, יהי $x \in \Jj(R)$ ונסמן $I = (x)$, אם $I \subset \Nil(R)$ אז סיימנו, לכן נניח ש־$I \not\subset \Nil(R)$.
	מהנתון קיים $y \in I$ כך ש־$y^2 = y$, ונסמן $y = x z$ ולכן,
	\[
		x z = x z x z
		\implies x z (1 - x z) = 0
	\]
	ונתון כי $1 - x z$ הפיך כאיבר ברדיקל ג'ייקובסון, אז נסיק $x z = 0$, ולכן בהכרח $x \in \Nil(R)$.
\end{proof}

\question{}
\subquestion{}
יהי $R$ חוג כך ש־$\Jj(R)$ סופי וכן שיש מספר אידאלים מקסימליים סופי. \\
נראה שלכל פולינום $f \in R[x]$ שאחד ממקדמיו הפיך יש כמות סופית של שורשים ב־$R$.
\begin{proof}
	נבחין כי אם $M_i \triangleleft R$ כלל האידאלים המקסימליים עבור $1 \le i \le n$ אז ממשפט השאריות הסיני,
	\[
		R / J(R)
		= R / \bigcap_{i = 1}^n M_i
		\cong \prod_{i = 1}^n (R / M_i)
	\]
	נגדיר את $\varphi : R \to \prod R / M_i$ ולכן $\ker \varphi = J(R)$.
	ל־$f$ יש מקדם הופכי ולכן גם ל־$f + M_i \in R[x] / M_i$ יש מקדם הפיך ולא אפס ולכן בפרט $f$ לא טריוויאלי במכפלה.
	נסמן $\deg f = k$ ולכן יש ל־$f$ לכל היותר $k$ שורשים בכל $R / M_i$ ובסך־הכול לכל היותר $k^n$ שורשים במכפלת השדות.
	נבחין כי אם $x \in \im \varphi$ אז יש לו לכל היותר $|J(R)|$ מקורות ממשפט האיזומורפיזם הראשון, אבל זהו מספר סופי.
\end{proof}

\subquestion{}
נראה ש־$R$ הוא חוג שבו $\Jj(R) = 0$ ויש לו כמות סופית של אידאלים מקסימליים אם ורק אם יש אוסף סופי של שדות $F_i$ ל־$1 \le i \le n$ כך ש־$R \cong \prod F_i$.
\begin{proof}
	הכיוון הראשון נובע מיידית ממשפט השאריות הסיני, נראה את הכיוון ההפוך. \\
	נסמן $\pi_j : \prod F_i \to F_j$ להיות ההטלה הסטנדרטית, וכן $\varphi : R \to \prod F_i$ איזומורפיזם.
	נגדיר לכל $1 \le i \le n$ את הפונקציה $p_i : R \to F_i$ על־ידי $p_i = \pi_i \circ \varphi$.
	זוהי הרכבה של הומומורפיזמים ולכן $M_i : = \ker p_i$ אידאל, וברור כי הוא מקסימלי וכי כל אידאל מקסימלי על־ידי $\varphi$ הוא מהצורה הזו, לכן מצאנו שיש מספר סופי של אידאלים סופיים.
	לבסוף נבחין כי גם,
	\[
		\Jj(R)
		= \bigcap M_i
		= \bigcap p_i^{-1}(\{ 0 \})
		= \{ 0 \}
	\]
	שכן נקבל את הדרישה שיהיה אפס בכל קורדינטה.
\end{proof}

\question{}
יהיו $I, J \triangleleft R$ אידאלים.

\subquestion{}
נראה שאם יש $n \in \NN$ שעבורו $I^n \subset J$ אז $\sqrt{I} \subset \sqrt{J}$.
\begin{proof}
	יהי $x \in I$, אז מהנתון $x^n \in J$ ולכן מהגדרה גם $x \in \sqrt{J}$.
	אם $y \in \sqrt{I}$ אז קיים $m \in \NN$ כך ש־$y^m \in I$ ולכן $y^{n m} \in J$.
\end{proof}

\subquestion{}
נראה שמתקיים,
\[
	\sqrt{\sqrt{I}}
	= \sqrt{I}
\]
\begin{proof}
	יהי $x \in \sqrt{\sqrt{I}}$, לכן קיים $y \in \sqrt{I}$ כך ש־$x^n = y$, ולכן קיים $z \in I$ כך ש־$y^m = z$, אבל בהתאם $x^{n m} \in I$ ולכן $\sqrt{\sqrt{I}} \subset \sqrt{I}$.
	הצד השני נובע ישירות מ־$I \subset \sqrt{I}$.
\end{proof}

\subquestion{}
נראה שמתקיים,
\[
	\sqrt{I J}
	= \sqrt{I \cap J}
	= \sqrt{I} \cap \sqrt{J}
\]
\begin{proof}
	נובע בצורה זהה לסעיף הקודם.
\end{proof}

\subquestion{}
נוכיח ש־$\sqrt{I} = (1)$ אם ורק אם $I = (1)$.
\begin{proof}
	\[
		\sqrt{I} = (1)
		\iff 1 \in \sqrt{I}
		\iff \exists x \in I \exists n \in \NN,\ 1^n = x
	\]
	אבל $1^n = 1$ לכל $n \in \NN$.
\end{proof}

\subquestion{}
נראה שמתקיים,
\[
	\sqrt{I + J}
	= \sqrt{\sqrt{I} + \sqrt{J}}
\]
\begin{proof}
	נובע מסגירות של חיבור אידאלים לחיבור וסעיף ב'.
\end{proof}

\subquestion{}
נוכיח שאם $I$ ראשוני אז לכל $n \in \NN$ מתקיים $\sqrt{I^n} = I$.
\begin{proof}
	אם $x \in \sqrt{I^n}$ אז $x^m \in I^n$ לאיזשהו $m \in \NN$.
	אבל במצב זה נקבל ש־$x^{m - 1} \in I$ או $x \in I$, ולאחר מספר סופי של מהלכים כאלה נקבל ש־$x \in I$.
	הכיוון השני דומה.
\end{proof}

\question{}
יהי $R$ חוג ו־$I \triangleleft R$, נראה שמתקיים,
\[
	\sqrt{I}
	= \bigcap_{P \in \Pp} P
\]
עבור $\Pp$ משפחת האידאלים הראשוניים ב־$R$ המכילים את $I$.
\begin{proof}
	יהי $P \in \Pp$ ונראה ש־$\sqrt{I} \subset P$.
	יהי $x \in \sqrt{I}$ כך ש־$x^n \in I$, אז $x^n \in P$ ולכן $x \in P$ או $x^{n - 1} \in P$, לאחר מספר סופי של צעדים נקבל שבהכרח $x \in P$.
	נסיק שמתקיים $\sqrt{I} \subset \bigcap_{P \in \Pp} P$.

	בכיוון ההפוך יהי $x \notin \sqrt{I}$ ולכן בפרט $x \notin I$, ונגדיר את $S$ להיות כל האידאלים המכילים את $I$ וזרים ל־$x, x^2, \ldots$.
	נבחין כי $I \in S$ ולכן נוכל להשתמש בלמה של צורן עם הכלה ושימוש בעובדה שאיחוד אידאלים הוא אידאל, ונסמן $P \in S$ איבר מקסימלי.
	מאותו שיקול של משפט קיום אידאל מקסימלי נוכל לקבל שגם $P$ הוא ראשוני ולכן מעיד על $x \notin \bigcap \Pp$.
\end{proof}

\question{}
יהיו $p_1, \ldots, p_k \in \ZZ$ ראשוניים שונים, ויהיו $n_1, \ldots, n_k \in \NN$ שונים מאפס.
נראה שמתקיים,
\[
	\sqrt{p_1^{n_1} \cdots p_k^{n_k} \ZZ}
	= p_1 \cdots p_k \ZZ
\]
\begin{proof}
	\[
		\sqrt{p_1^{n_1} \cdots p_k^{n_k} \ZZ}
		= \sqrt{p_1^{n_1} \ZZ} \cap \cdots \cap \sqrt{p_k^{n_k} \ZZ}
		= p_1 \ZZ \cap \cdots \cap p_k \ZZ
		= p_1 \cdots p_k \ZZ
	\]
\end{proof}

\question{}
יהי $f = \sum_{i = 0}^n a_i x^i \in R[x]$.

\subquestion{}
נראה ש־$f$ נילפוטנטי אם ורק אם $a_0, \ldots, a_n$ נילפוטנטיים.
\begin{proof}
	אם $a_0, \ldots, a_n$ נילפוטנטיים אז יש טבעי יחיד $k \in \NN$ עבורו $a_i^k = 0$ לכל $i$, עבורו $f^k = 0$ בלבד.

	נניח ש־$f$ נילפוטנטי ולכן קיים $k$ עבורו $f^k = 0$, אבל גם,
	\[
		f^k
		= {\left(\sum_{i = 0}^n a_i x^i\right)}^k
		= \sum_{\substack{(i_0, \ldots, i_n) \in \NN^n, \\ i_0 + \cdots + i_n = k}} a_0^{i_0} \cdots a_n^{i_n} x^{i_1 + 2 i_2 + \cdots + n i_n}
	\]
	ולכן $a_i^k = 0$ לכל $i$.
\end{proof}

\subquestion{}
נראה ש־$f$ מחלק אפס אם ורק אם קיים $0 \ne a \in R$ כך ש־$a f = 0$.
\begin{proof}
	יהי $0 \ne g \in R[x]$ ממעלה מינימלית כך ש־$f g = 0$, ונניח בשלילה ש־$\deg g > 0$.
	נסמן $g = \sum_{i = 0}^m b_i x^i$ ולכן $f g = a_0 b_0 + \cdots + a_n b_m x^{n + m}$, ובפרט $a_n b_m = 0$ בלבד.
	נובע אם כך ש־$a_n g$ הוא פולינום כך ש־$\deg a_n g < \deg g$ וגם $b_n g f = b_n (f g) = 0$, וזו סתירה, לכן $\deg g = 0$.
\end{proof}

\question{}
יהיו $I = (4, x), M = (2, x) \triangleleft \ZZ[x]$ אידאלים.

\subquestion{}
נראה ש־$M$ מקסימלי.
\begin{proof}
	$\ZZ[x] / (2, x) \cong (\ZZ[x] / (x)) / (2) = \ZZ / (2) \cong \FF_2$, וזהו שדה ולכן $M$ מקסימלי.
\end{proof}

\subquestion{}
נוכיח ש־$M = \sqrt{I}$.
\begin{proof}
	נבחין כי $2^2 = 4, x^1 = x$ ולכן $M \subset \sqrt{I}$.
	מתקיים גם $3 \notin \sqrt{I}$ שכן $3 \notin I$ ולא קיים $n^k = 3$, לכן $M \subset \sqrt{I} \subsetneq \ZZ[x]$ ולכן $M = \sqrt{I}$ בלבד.
\end{proof}

\subquestion{}
נראה שלא קיים $n \in \NN$ כך ש־$M^n = I$.
\begin{proof}
	עבור $n = 1$ הטענה ברורה, ועבור $n > 1$ נקבל ש־$x \notin M^n$ ובפרט $M^n \ne I$.
\end{proof}

\question{}
יהי $F$ שדה, לכל אידאל נבדוק אם הוא ראשוני או מקסימלי.

\subquestion{}
נראה ש־$(x - y - 1) \triangleleft F[x, y]$ הוא ראשוני.
\begin{proof}
	נגדיר את ההומומורפיזם $\varphi : F[x, y] \to F[x]$ על־ידי $\varphi(f) = f(x, x - 1)$.
	לכל $f \in F[x]$ מתקיים $\varphi(f) = f$ ולכן $\varphi$ הוא על. \\
	נבחין כי $\varphi(x - y - 1) = x - (x - 1) - 1 = 0$ ולכן $(x - y - 1) \subset \ker \varphi$, נרצה להראות גם הכלה הפוכה.
	יהי $f \in \ker \varphi$ ונסמן $f = \sum_{i, j} c_{i, j} x^i y^j$ אז,
	\begin{align*}
		f + (x - y - 1)
		& = \sum_{i, j} c_{i, j} x^i y^j + (x - y - 1) \\
		& = \sum_{i, j} c_{i, j} x^i {(x - 1)}^j + (x - y - 1) \\
		& = f(x, x - 1) + (x - y - 1)
	\end{align*}
	ולכן $\ker \varphi = (x - y - 1)$ בדיוק ונובע ש־$F[x] / (x - y - 1) = F[x]$, זהו תחום שלמות שאיננו שדה ולכן $(x - y - 1)$ ראשוני אבל לא מקסימלי.
\end{proof}

\subquestion{}
נראה ש־$(x, y^2 - 1) \triangleleft F[x, y]$ הוא מקסימלי.
\begin{proof}
	נבחין כי $F[x, y] / (x, y^2 - 1) \cong (F[x, y] / x) / (y^2 - 1) \cong F[y] / (y^2 - 1) \cong F[i]$ ולכן מספיק להראות שכל איבר ב־$F[i]$ הפיך.
	יהי $x + i y \in F[x]$ אז $\frac{(x + i y)(x - i y)}{x^2 + y^2} = 1$ וכן $\frac{x}{x^2 + y^2}, \frac{y}{x^2 + y^2} \in F$ ולכן קיים הופכי וזהו שדה, כלומר $(x, y^2 - 1)$ אידאל מקסימלי.
\end{proof}

\subquestion{}
נראה ש־$(x - y + 1, y - 3) \triangleleft F[x, y]$ הוא אידאל מקסימלי.
\begin{proof}
	נבחין כי $(x - y + 1, y - 3) = (x - 2, y - 3)$ שכן $x - y + 1 - (y - 3) = x - 2$.
	לכן,
	\[
		F[x, y] / (x - 2, y - 3)
		\cong (F[x, y] / (x - 2)) / (y - 3)
		\cong \FF_2[y] / (y - 3)
		= \FF_2[y] / (y - 1)
		\cong \FF_2
	\]
	זהו שדה ולכן $(x - y + 1, y - 3)$ הוא אידאל מקסימלי.
\end{proof}

\subquestion{}
נראה ש־$(2 + x + z, y - x z - x^2) \triangleleft F[x, y, z]$ הוא אידאל ראשוני.
\begin{proof}
	מתקיים $y - x z - x^2 + x (2 + x + z) = y + 2 x$ ולכן $(2 + x + z, y - x z - x^2) = (2 + x + z, y + 2 x)$.
	נגדיר את ההומומורפיזם $\varphi : R[x, y, z] \to R[x]$ על־ידי $\varphi(f) = f(x, -2 y, -x - 2)$. 
	ברור כי $\im \varphi = F[x]$ ונראה ש־$\ker \varphi = (2 + x + z, y + 2 x) = I$.
	מתקיים $\varphi(2 + x + z) = 0, \varphi(y + 2 x) = 0$ ולכן $I \subset \ker \varphi$ ונותר להראות את הכיוון ההפוך.
	יהי $f \in \ker \varphi$ ונסמן $f = \sum_{i, j, k} c_{i, j, k} x^i y^j z^k$ אז,
	\[
		f + I
		= \sum_{i, j, k} c_{i, j, k} x^i y^j z^k + I
		= \sum_{i, j, k} -2 c_{i, j, k} x^{i + 1} y^{j - 1} (-x - 2) z^{k - 1} + I
		= f(x, -2x, -x - 2) + I
		= 0
	\]
	לכן ממשפט האיזומורפיזם הראשון $F[x, y, z] / I \cong F[x]$ ולכן $I$ אידאל ראשוני שאינו מקסימלי.
\end{proof}

\end{document}
